\section{The Recursion}

\textbf{[STIPULATION]} The scale structure of the theory is generated by the map
\[
  \xi_{n+1} = \xi_n\,\bigl(1 - 100\,\xi_n\bigr)
\]
starting from $\xi_0 = 4/30000$.

This is a damped feedback: the factor $(1-100\xi_n)$ is less than one, the
sequence falls monotonically. In communications engineering terms it is a loop
with gain below one -- it contracts, it does not oscillate. The gain is not
adjustable: it depends solely on $\xi$, which is stipulated. There is no tuning
knob in this loop.

\section{The Non-Closure Theorem}

\textbf{[K]} The recursion has exactly one fixed point, namely zero. Every
value $c \in (0,1)$ that the sequence assumes is a transit point, not a target.

The proof is elementary: from $\xi = \xi(1-100\xi)$ it follows that $100\xi^2 = 0$,
hence $\xi = 0$. And because the sequence is strictly monotonically decreasing
but stays positive, it will not revisit any of its values in finitely many steps
and will not reach zero in any finite step.

\textbf{[K]} From this follows a statement that is easy to overlook and yet
load-bearing: \emph{from the recursion alone, no individual numerical value can
be singled out}. For every target $c$ there exists a depth $k^*(c)$ at which
the sequence comes close to it -- equally so for every $c$, with identical
parameter status. One who tries to justify a number by claiming the recursion
passes through it has justified nothing.

\textbf{[K]} This is a result, not a gap. It closes an entire class of derivation
attempts: the question of whether a specific constant can be obtained from the
recursion depth is \emph{answered}, and the answer is no. What remains open is
only whether there is an anchor outside the recursion.

\section{The Frozen Winding}

\textbf{[B]} Holding $\xi$ fixed instead of letting it run, the winding closes
after exactly
\[
  \frac{1}{100\,\xi_0} = 75
\]
circuits. The rotation number is then $74/75$, and the missing piece is the
closure comma of A040.

This is the \emph{rolled-up}, compact reading: one scale, one closed winding,
an integer number of circuits.

\section{The Unrolled Reading: A Logarithmic Spiral}

\textbf{[K]} Letting $\xi$ run, the sequence decays asymptotically as
\[
  \xi_n \approx \frac{1}{100\,(n+75)} ,
\]
and the cumulative deficit becomes a harmonic sum, hence logarithmic:
\[
  D(k) \approx \ln\!\Bigl(1 + \frac{k}{75}\Bigr).
\]

\textbf{[K]} Setting $\rho = k + 75$ and $\beta = 2\pi D$, we get
\[
  \rho = 75\,e^{\beta/2\pi} ,
\]
a logarithmic spiral with self-similarity ratio $e$.

Two points, both important.

The ratio $e$ is \emph{not an additional parameter}. It is the base of the
natural logarithm and comes from the harmonic sum -- it is not chosen, it falls
out. The pole at $n = -75 = -1/(100\xi_0)$ is likewise not fitted but fixed
by $\xi_0$.

The scale invariance is \emph{continuous}, not discrete log-periodic. There are
no distinguished steps, no preferred multiples.

\textbf{[B]} Rolled up and unrolled are two readings of the same object, not
two objects: the closed $75$-winding is the compact single-scale view, the
spiral the unrolled across-scales view.

\section{Why the Recursion Does Not Belong on a Single Axis}

\textbf{[B]} The formulation that time rolls up suggests that the time axis is
distinguished. It is not. Through $\tilde T\cdot m = 1$, every statement about
time has a mirror image about mass: if time advances by $K_\text{frak}$ per
step, mass advances by $1/K_\text{frak}$. The deficit is the same $100\xi_n$
on both sides.

\textbf{[B]} The factor $100$ therefore does not migrate between the axes -- it
is mirrored. Both projections lead to the same ratio $e$. The assignment to time
is a convention, not a distinction.

\section{Verification Scripts}

\begin{itemize}
  \item Fixed-point theorem and non-closure: unique fixed point $0$, monotonicity,
    and the construction of $k^*(c)$ for arbitrary targets -- as evidence that
    no value is distinguished:
    \texttt{python/A\_Serie\_Skripte/a050\_nichtschliessung.py}
  \item Spiral ratio: numerical determination of the ratio from the cumulative
    deficit, comparison with $e$, and check for continuous rather than
    log-periodic scale invariance:
    \texttt{python/A\_Serie\_Skripte/a050\_spirale\_e.py}
\end{itemize}

\section{Open Questions}

\textbf{First, the recursion depth is not fixed.} The theory gives no rule
at which $k$ a physical quantity is to be read off. By the non-closure theorem
it cannot -- this is not a negligence, but proved.

\textbf{Second, this implies a hard limit for all fine values.} Every quantity
whose justification would amount to the recursion reaching it is thereby
\emph{not} justified. This applies notably to the fine value of the refraction
strength $\delta$; it remains open and is booked in A230.

\textbf{Third, the recursion uses the factor $100$, which is derived in A040}
(cumulative RG run, consistency condition; A040). It is used here, not
re-derived. The non-closure holds independently of this -- it follows from
the form of the map, not from the numerical value -- but the concrete number
$75$ does not.

\section{Supersedes}

This document subsumes: Dok.~203 (R-operator), Dok.~275 ($\xi$-recursion),
Dok.~295 (time as logarithmic spiral). Incorporated: P37 (three roles of
$\varphi$ separated), P35 (triviality of the power representation).

\section{Use of AI Tools}

This document was created with the assistance of AI tools.
Responsibility for content and errors rests entirely with the author.
Full wording of the rules in A010.
